\documentclass[conference]{IEEEtran}
\IEEEoverridecommandlockouts

\usepackage{cite}
\usepackage{amsmath,amssymb,amsfonts}
\usepackage{graphicx}
\usepackage{booktabs}
\usepackage{multirow}
\usepackage{textcomp}
\usepackage{xcolor}
\usepackage{url}

% ---------------------------------------------------------------------------
% Drafting helpers -- remove before camera-ready submission.
%   \tbd{...}   : number or text that depends on experiments still running
%   \figslot    : includes the figure if the file exists, otherwise a framed placeholder
% ---------------------------------------------------------------------------
\newcommand{\tbd}[1]{\textcolor{red}{[#1]}}
\newcommand{\figslot}[3]{%
  \IfFileExists{#1}{\includegraphics[width=#2]{#1}}%
  {\fbox{\parbox[c][3cm][c]{#2}{\centering\footnotesize\textcolor{red}{FIGURE: #3}}}}}

\begin{document}

% Title kept as submitted. The text (abstract, Sec. III-A) states that SE and ResNet-D
% are evaluated separately, since no pretrained SE-ResNet-50-D exists.
\title{Stabilized Reinforcement Learning for Image Captioning using SE-ResNet-D}

\author{
\IEEEauthorblockN{Manas Mehta}
\IEEEauthorblockA{\textit{CSE, GLS University}\\
Ahmedabad, India \\
manasmehta1110@gmail.com}
\and
\IEEEauthorblockN{Vishvesh Sharma}
\IEEEauthorblockA{\textit{CSE, GLS University}\\
Ahmedabad, India \\
sharmavishvesh3@gmail.com}
\and
\IEEEauthorblockN{Dr. Madhuri Chopade}
\IEEEauthorblockA{\textit{CSE, GLS University}\\
Ahmedabad, India \\
madhuri.chopade@glsuniversity.ac.in}
}

\maketitle

\begin{abstract}
Self-Critical Sequence Training (SCST) lets an image captioning model optimize
its evaluation metric directly, but the switch from cross-entropy (CE) training
to reinforcement learning can destabilize the model and push it toward
repetitive, metric-seeking captions. We study this transition in a lightweight
attention-based CNN-LSTM captioner trained on a single 4\,GB consumer GPU.
For the encoder, we compare ImageNet-pretrained ResNet-50 models with and
without Squeeze-and-Excitation (SE) blocks and the ResNet-D stem and
downsampling path, all pretrained with the same recipe. For training, we analyze a mixed objective that weights the SCST loss by 0.8 and
the CE loss by 0.2, combined with normalized advantages and an optimizer reset
at the CE-to-RL switch, and we ablate each component. On the Karpathy test
split of MS-COCO, the ResNet-D encoder gives the best CE model, with 31.3
BLEU-4 and 98.2 CIDEr, while SE blocks give no measurable gain. Beyond the
standard metrics, we measure reward hacking directly. Pure SCST raises CIDEr
but ends 77.9\% of its test captions with a function word such as ``a'' or
``with''. The mixed objective keeps this rate at zero for a cost of 1.8 CIDEr,
and normalizing the advantage makes the switch to reinforcement learning
smoother.
\end{abstract}

\begin{IEEEkeywords}
image captioning, self-critical sequence training, reinforcement learning,
squeeze-and-excitation, ResNet-D, attention, MS-COCO
\end{IEEEkeywords}

% ===========================================================================
\section{Introduction}
% ===========================================================================
Image captioning requires a model to recognize the objects in a scene, relate
them to one another, and express the result as a fluent sentence. The
encoder-decoder paradigm introduced by Show and Tell~\cite{vinyals2015show}
and extended with visual attention in Show, Attend and
Tell~\cite{xu2015show} remains a standard design: a convolutional neural
network (CNN) encodes the image into a grid of feature vectors, and a
long short-term memory (LSTM) decoder attends to that grid while generating one
word at a time.

Captioning models are usually trained with a cross-entropy (CE) loss under
teacher forcing, where the decoder always receives the ground-truth previous
word. At inference time, it must instead condition on its own predictions, a
mismatch known as exposure bias~\cite{ranzato2016sequence,bengio2015scheduled}.
CE training also optimizes word-level likelihood, whereas captions are judged
by sequence-level metrics such as BLEU~\cite{papineni2002bleu} and
CIDEr~\cite{vedantam2015cider}. Self-Critical Sequence Training
(SCST)~\cite{rennie2017self} addresses both issues by optimizing the metric
directly with the REINFORCE estimator, using the model's own greedy caption as
a baseline. SCST has problems of its own. The model can learn to exploit the
reward by producing generic or repetitive captions that score well on CIDEr
(reward hacking), and the abrupt change of objective at the CE-to-RL switch can
destabilize a model trained under a different loss.

We study both stages of this pipeline with limited compute. Our contributions
are:
\begin{itemize}
\item A controlled comparison of ResNet-50~\cite{he2016deep},
SE-ResNet-50~\cite{hu2018squeeze}, and ResNet-50-D~\cite{he2019bag} encoders,
pretrained with the same recipe~\cite{wightman2021resnet} and paired with the
same decoder.
\item An ablation of a stabilized SCST setup: a mixed 80/20 SCST/CE
objective~\cite{paulus2018deep}, normalized advantages, and an optimizer reset
at the CE-to-RL switch.
\item A direct measurement of reward hacking through vocabulary usage, caption
uniqueness, $n$-gram diversity, and repetition.
\end{itemize}
Sample captions and attention maps are shown in Fig.~\ref{fig:attention}.

\begin{figure}[t]
\centering
\figslot{figs/attention_maps.pdf}{\columnwidth}{attention maps}
\caption{Attention of the mixed SCST model on randomly drawn Karpathy test
images, with the greedy caption above each image. Each map brightens the
regions weighted by $\alpha_t$ when the word above it was generated.}
\label{fig:attention}
\end{figure}

% ===========================================================================
\section{Related Work}
% ===========================================================================
\textit{Encoder-decoder captioning.} Early neural captioners paired a CNN
image encoder with a recurrent language model~\cite{vinyals2015show,karpathy2015deep,mao2015deep}.
Show, Attend and Tell~\cite{xu2015show} added soft attention over a spatial
feature grid, so that the decoder can focus on different image regions at each
step. Anderson et al.~\cite{anderson2018bottom} replaced the uniform grid with
object-level features from a detector (bottom-up attention). Later models
refined the attention mechanism itself~\cite{huang2019attention} and moved to
fully transformer-based encoders and decoders~\cite{cornia2020meshed}.
These models score much higher (Table~\ref{tab:main}), but they rely on object
detectors or large-scale pretraining and need far more compute than we target
here.

\textit{Encoder architectures.} Squeeze-and-Excitation (SE)
blocks~\cite{hu2018squeeze} recalibrate channels with a gate computed from
globally pooled features, and the ResNet-D variant of He et
al.~\cite{he2019bag} changes the stem and the downsampling shortcut
(Section~\ref{sec:encoder}). Both improve ImageNet accuracy; we test whether
they also improve captioning when the decoder is held fixed.

\textit{Sequence-level training.} Ranzato et al.~\cite{ranzato2016sequence}
used REINFORCE to optimize sequence-level metrics directly. SCST~\cite{rennie2017self}
reduced the variance of this estimator by using the reward of the model's own
greedy decoding as the baseline, and it has become the standard final stage of
captioning pipelines~\cite{anderson2018bottom,huang2019attention,cornia2020meshed}.
Luo~\cite{luo2020better} replaced the greedy baseline with the mean reward of
multiple samples. Paulus et al.~\cite{paulus2018deep} combined maximum-likelihood
and RL losses to keep generated summaries readable, which is the idea behind
the mixed objective we study.

% ===========================================================================
\section{Method}
% ===========================================================================
Fig.~\ref{fig:arch} gives an overview of the model. A frozen, ImageNet-pretrained
CNN encodes the image into a $7\times7$ grid of feature vectors, and an
attention-based LSTM decoder generates the caption.

\begin{figure*}[t]
\centering
\figslot{figs/architecture.pdf}{0.85\textwidth}{model overview}
\caption{Model overview. A frozen ImageNet-pretrained encoder produces a
$7\times7\times2048$ feature map, which is projected to 49 vectors of dimension
512. At each step, the LSTM decoder attends over these vectors and generates
the next word.}
\label{fig:arch}
\end{figure*}

\subsection{Encoder}
\label{sec:encoder}
We compare three encoders that share the ResNet-50
backbone~\cite{he2016deep}:
\begin{itemize}
\item \textbf{ResNet-50}, the baseline.
\item \textbf{SE-ResNet-50}, which adds an SE block~\cite{hu2018squeeze} to
every bottleneck. For a feature map $\mathbf{U}\in\mathbb{R}^{C\times H\times W}$,
the block computes a channel descriptor by global average pooling, passes it
through a two-layer bottleneck with reduction ratio $r=16$, and rescales each
channel:
\begin{align}
z_c &= \frac{1}{HW}\sum_{i=1}^{H}\sum_{j=1}^{W} u_c(i,j), \label{eq:squeeze}\\
\mathbf{s} &= \sigma\!\left(\mathbf{W}_2\,\delta(\mathbf{W}_1\mathbf{z})\right), \qquad
\tilde{\mathbf{u}}_c = s_c\,\mathbf{u}_c, \label{eq:excite}
\end{align}
where $\delta$ is the ReLU, $\sigma$ the sigmoid,
$\mathbf{W}_1\in\mathbb{R}^{C/r\times C}$ and
$\mathbf{W}_2\in\mathbb{R}^{C\times C/r}$.
\item \textbf{ResNet-50-D}~\cite{he2019bag}, which replaces the $7\times7$,
stride-2 stem convolution with three $3\times3$ convolutions and inserts a
$2\times2$ average pooling before the $1\times1$ projection in the
downsampling shortcut (Fig.~\ref{fig:stem}). In both our baseline and
ResNet-50-D, the stride of each downsampling bottleneck sits in its
$3\times3$ convolution, so these two changes are the only difference
between them.
\end{itemize}
All three encoders are pretrained on ImageNet-1k with the same recipe (the
``A1'' procedure of~\cite{wightman2021resnet}), which keeps pretraining
differences out of the comparison. We know of no publicly available
SE-ResNet-50-D pretrained with this recipe, so we evaluate the two
modifications separately.

The encoders are kept frozen. A $256\times256$ input yields an
$8\times8\times2048$ feature map, which we average-pool to $7\times7$ and
reshape into $L=49$ vectors $\mathbf{v}_k\in\mathbb{R}^{2048}$. A trainable
linear layer projects these to $\mathbf{a}_k = \mathbf{W}_p\mathbf{v}_k +
\mathbf{b}_p\in\mathbb{R}^{512}$.

\begin{figure}[t]
\centering
\figslot{figs/stem.pdf}{0.85\columnwidth}{ResNet-D modifications}
\caption{The two ResNet-D modifications. (a) The $7\times7$ stem
convolution is replaced by three $3\times3$ convolutions. (b) In the
downsampling shortcut, a stride-2 $1\times1$ convolution, which skips three of
every four positions, is replaced by $2\times2$ average pooling followed by a
stride-1 $1\times1$ convolution.}
\label{fig:stem}
\end{figure}

\subsection{Attention-Based LSTM Decoder}
The decoder follows the soft-attention model of Xu et al.~\cite{xu2015show}.
Its hidden and memory states are initialized from the mean feature,
$\mathbf{h}_0 = \mathbf{W}_{h}\bar{\mathbf{a}}$ and
$\mathbf{c}_0 = \mathbf{W}_{c}\bar{\mathbf{a}}$ with
$\bar{\mathbf{a}} = \frac{1}{L}\sum_k \mathbf{a}_k$. At step $t$, an attention
score is computed for every location from the previous hidden state, normalized
with a softmax, and used to form a context vector:
\begin{align}
e_{t,k} &= \mathbf{w}^{\top}\,\mathrm{ReLU}\!\left(\mathbf{W}_{a}\mathbf{a}_k + \mathbf{W}_{h}'\mathbf{h}_{t-1}\right), \label{eq:energy}\\
\alpha_{t,k} &= \frac{\exp(e_{t,k})}{\sum_{j=1}^{L}\exp(e_{t,j})}, \qquad
\hat{\mathbf{z}}_t = \sum_{k=1}^{L}\alpha_{t,k}\,\mathbf{a}_k. \label{eq:context}
\end{align}
As in~\cite{xu2015show}, a gate computed from the previous hidden state scales
the context vector element-wise before it enters the LSTM:
\begin{equation}
\boldsymbol{\beta}_t = \sigma\!\left(\mathbf{W}_{\beta}\mathbf{h}_{t-1}\right), \qquad
\mathbf{x}_t = \left[\mathbf{E}y_{t-1};\ \boldsymbol{\beta}_t\odot\hat{\mathbf{z}}_t\right],
\label{eq:gate}
\end{equation}
where $\mathbf{E}y_{t-1}$ is the embedding of the previous word. The gate lets
the decoder reduce its reliance on visual context, for example when it
generates function words such as ``the'' or ``of''. The LSTM state is updated
from $\mathbf{x}_t$, and the next-word distribution is
$p(y_t\mid y_{<t}, I) = \mathrm{softmax}(\mathbf{W}_o\,\mathrm{Dropout}(\mathbf{h}_t))$.
The embedding, hidden, and attention dimensions are all 512.

\subsection{Training Objectives}
\textit{Cross-entropy.} The decoder is first trained with teacher forcing to
minimize the token-level negative log-likelihood of the ground-truth caption
$y^{*}$:
\begin{equation}
L_{\mathrm{CE}} = -\frac{1}{T}\sum_{t=1}^{T}\log p_\theta\!\left(y^{*}_t\mid y^{*}_{<t}, I\right).
\label{eq:ce}
\end{equation}

\textit{Self-critical sequence training.} For each image, we sample a caption
$w^{s}$ from $p_\theta$ and decode a greedy caption $w^{g}$. Each caption is
scored with the CIDEr-D reward $r(\cdot)$, computed against all reference
captions of the image, with document frequencies precomputed over the full
training set. The greedy caption's reward is the
baseline~\cite{rennie2017self}:
\begin{equation}
L_{\mathrm{SCST}} = -\frac{1}{B}\sum_{i=1}^{B}\hat{A}_i\,\log p_\theta\!\left(w^{s}_i\right),
\label{eq:scst}
\end{equation}
where $\log p_\theta(w^{s}_i)$ sums the token log-probabilities up to and
including the first end-of-sentence token, and $B$ is the batch size.

\textit{Normalized advantage.} Instead of the raw advantage
$A_i = r(w^{s}_i) - r(w^{g}_i)$, we use its batch-standardized form
\begin{equation}
\hat{A}_i = \frac{A_i - \mu_A}{\sigma_A + \epsilon},
\label{eq:adv}
\end{equation}
where $\mu_A$ and $\sigma_A$ are the mean and standard deviation of $A$ within
the batch and $\epsilon=10^{-8}$. This keeps the scale of the policy gradient
constant as rewards grow during training.

\textit{Mixed objective.} Pure SCST removes the likelihood term that keeps
captions close to human phrasing. Following~\cite{paulus2018deep}, we keep a
fraction of it:
\begin{equation}
L = \lambda\,L_{\mathrm{SCST}} + (1-\lambda)\,L_{\mathrm{CE}}, \qquad \lambda = 0.8,
\label{eq:mixed}
\end{equation}
where the CE term uses one randomly chosen reference caption per image.

\textit{Optimizer reset.} The Adam moment estimates accumulated during CE
training are calibrated to CE gradients. At the CE-to-SCST switch, we discard
them and start a fresh Adam optimizer with the lower SCST learning rate.

% ===========================================================================
\section{Experimental Setup}
% ===========================================================================
\subsection{Dataset and Preprocessing}
We use MS-COCO~\cite{lin2014microsoft} with the Karpathy
split~\cite{karpathy2015deep}: 113{,}287 training images (including the
``restval'' images), 5{,}000 validation images, and 5{,}000 test images, each
with at least five reference captions. Captions are lower-cased and tokenized,
and words that occur five times or fewer are mapped to an unknown token, giving
a vocabulary of 9{,}490 tokens including the start, end, unknown, and padding
tokens. Captions longer than 50 words are discarded for training. Images are
resized to $256\times256$ and normalized with the ImageNet mean and standard
deviation. Because the encoders are frozen, their features are computed once
and cached, and no data augmentation is applied.

\subsection{Training Details}
\label{sec:training}
All models are implemented in PyTorch and trained on a single NVIDIA GeForce
RTX 3050 Laptop GPU with 4\,GB of memory. Pretrained encoder weights are taken
from the \texttt{timm} library. We use Adam with weight decay $10^{-4}$,
element-wise gradient clipping at 2.0, and dropout of 0.5 before the output
layer.

\textit{CE stage.} The decoder and projection layer are trained from scratch
with CE at a learning rate of $4\times10^{-4}$ for up to 20 epochs. Each batch
holds 16 images with all of their
reference captions, so every (image, caption) pair is seen once per epoch. The
learning rate is halved when the validation CIDEr fails to improve for two
consecutive epochs, and training stops after three epochs without improvement.

\textit{SCST stage.} All SCST experiments use the SE-ResNet-50 encoder, the one
closest to our original design. We fixed this choice before the encoder
comparison was complete. Starting from the best SE-ResNet-50 CE checkpoint, the
model is trained with the objective under study for 10 epochs, with one step
per batch of 32 images. Sampled and greedy captions are limited to 20
words. We set the SCST learning rate to $5\times10^{-5}$, as
in~\cite{rennie2017self}. In a 300-step pilot run, the $10^{-6}$ used in our
original experiments left the validation CIDEr unchanged.

\textit{Model selection.} For every run, the checkpoint with the highest
validation CIDEr under greedy decoding is selected. The test split is used only
for the final evaluation of the selected checkpoints.

\subsection{Evaluation Protocol}
Captions are generated with beam search (beam size 7). Beam scores are
normalized by caption length raised to $\alpha=0.7$, the end token is blocked
for the first five words, and any word that would repeat an existing 4-gram is
blocked. We also report greedy decoding and plain beam search (beam size 3, no
length penalty or blocking) to show how much decoding choices contribute
(Section~\ref{sec:decoding}).

We compute BLEU-1 to BLEU-4~\cite{papineni2002bleu},
METEOR~\cite{banerjee2005meteor}, ROUGE-L~\cite{lin2004rouge}, and
CIDEr-D~\cite{vedantam2015cider} with the official COCO caption evaluation
code~\cite{chen2015microsoft}, using the Penn Treebank tokenizer and all
reference captions of each test image. Following common practice, scores are
reported multiplied by 100. Captioning has no single correct output, so
classification accuracy does not apply. BLEU measures $n$-gram precision
against the references, ROUGE-L is recall-oriented, METEOR combines precision
and recall, and CIDEr-D weights $n$-grams by TF-IDF.

To quantify reward hacking, we also report the number of distinct words used,
the share of unique captions, corpus-level distinct-2 (distinct bigrams over
all bigrams), the share of captions that repeat a content word, and the share
that end in a function word such as ``a'', ``with'', or ``of''.

% ===========================================================================
\section{Results and Analysis}
% ===========================================================================
\subsection{Main Results}
Table~\ref{tab:main} compares our models with published results on the
Karpathy test split. With a frozen ImageNet-pretrained SE-ResNet-50 or
ResNet-50-D encoder, our CE models reach 31.3 BLEU-4, above the CNN-LSTM
models of Karpathy and Fei-Fei (23.0) and of Xu et al. (24.3 and 25.0), and
equal to the Att2in model of Rennie et al.,
which uses a larger ResNet-101 encoder and SCST. Their CIDEr of 96.5 to 98.2 is
below Att2in's 101.3. Models built on object-detector features remain 22 to 33
CIDEr points higher, a gap we expect given the difference in features and
compute.

% VERIFY every literature number below against the original paper before submission.
% "--" = not reported. Soft/Hard attention METEOR were computed with an older
% METEOR version; see [xu2015show].
\begin{table*}[t]
\caption{Comparison on the MS-COCO Karpathy test split (scores $\times100$).
B-$n$: BLEU-$n$, M: METEOR, R: ROUGE-L, C: CIDEr-D. Models above the line use
CNN grid features and an RNN decoder; models below it use detector features
and/or transformers. For our models, ``(ImageNet)'' marks a frozen encoder
pretrained on ImageNet-1k, and ``(COCO only)'' an encoder trained from scratch
on MS-COCO without ImageNet pretraining.}
\label{tab:main}
\centering
\begin{tabular}{@{}llcccccc@{}}
\toprule
Model & Encoder & B-1 & B-4 & M & R & C & Training \\
\midrule
Deep VS~\cite{karpathy2015deep} & VGG-16 & 62.5 & 23.0 & 19.5 & -- & 66.0 & CE \\
Soft attention~\cite{xu2015show} & VGG-19 & 70.7 & 24.3 & 23.9 & -- & -- & CE \\
Hard attention~\cite{xu2015show} & VGG-19 & 71.8 & 25.0 & 23.0 & -- & -- & CE \\
Att2in~\cite{rennie2017self} & ResNet-101 & -- & 31.3 & 26.0 & 54.3 & 101.3 & SCST \\
\midrule
Up-Down~\cite{anderson2018bottom} & Faster R-CNN & 79.8 & 36.3 & 27.7 & 56.9 & 120.1 & SCST \\
AoANet~\cite{huang2019attention} & Faster R-CNN & 80.2 & 38.9 & 29.2 & 58.8 & 129.8 & SCST \\
$\mathcal{M}^2$ Transformer~\cite{cornia2020meshed} & Faster R-CNN & 80.8 & 39.1 & 29.2 & 58.6 & 131.2 & SCST \\
\midrule
Ours, from scratch$^\dagger$ (Sec.~\ref{sec:original}) & SE-ResNet-50 (COCO only) & 62.7 & 22.1 & 19.3 & 45.9 & 61.0 & mixed SCST \\
Ours, CE & SE-ResNet-50 (ImageNet) & 72.0 & 31.3 & 25.0 & 53.3 & 96.5 & CE \\
Ours, CE & ResNet-50-D (ImageNet) & 71.9 & 31.3 & 25.4 & 53.5 & 98.2 & CE \\
Ours, mixed SCST & SE-ResNet-50 (ImageNet) & 72.7 & 28.6 & 24.0 & 52.8 & 96.8 & 0.8 SCST + 0.2 CE \\
\bottomrule
\multicolumn{8}{@{}l}{\footnotesize $^\dagger$Deep $3\times3$ stem; no ImageNet pretraining, stem and first stage left at random initialization,}\\
\multicolumn{8}{@{}l}{\footnotesize \phantom{$^\dagger$}and SCST learning rate $10^{-6}$. Not comparable to the encoder ablation (Table~\ref{tab:encoder}).}
\end{tabular}
\end{table*}

\subsection{Encoder Ablation}
Table~\ref{tab:encoder} compares the three encoders under the CE stage, with the
decoder, data, and training budget held fixed. ResNet-50-D gives the best
CIDEr, 1.5 points above the ResNet-50 baseline, and the best METEOR and
ROUGE-L. Its lead holds under greedy decoding on the test split (90.6 versus
89.6 CIDEr) and on the validation split (90.2 versus 88.4). SE-ResNet-50 is
within 0.2 CIDEr of the baseline and slightly ahead on BLEU-4, so we see no
clear benefit from SE blocks when the encoder is frozen. The ResNet-D gain
holds across both splits and both decoding methods, but each encoder was
trained once, so we have not measured how much it varies between runs.

\begin{table}[t]
\caption{Encoder ablation after CE training (Karpathy test, scores $\times100$).
All encoders are frozen and pretrained on ImageNet-1k with the same recipe.}
\label{tab:encoder}
\centering
\setlength{\tabcolsep}{4pt}
\begin{tabular}{@{}lccccc@{}}
\toprule
Encoder & B-1 & B-4 & M & R & C \\
\midrule
ResNet-50 & 71.6 & 30.7 & 25.2 & 53.3 & 96.7 \\
SE-ResNet-50 & \textbf{72.0} & \textbf{31.3} & 25.0 & 53.3 & 96.5 \\
ResNet-50-D & 71.9 & \textbf{31.3} & \textbf{25.4} & \textbf{53.5} & \textbf{98.2} \\
\bottomrule
\end{tabular}
\end{table}

\subsection{Training-Objective Ablation}
Table~\ref{tab:loss} compares the training objectives, starting from the same
CE checkpoint with the SE-ResNet-50 encoder.

\textit{SCST versus CE.} Pure and mixed SCST both remove most repetition: the
share of captions with a repeated content word falls from 13.4\% to 2.5\%
(mixed) and 1.3\% (pure). Their effect on the scores depends on decoding. With
greedy decoding, mixed SCST raises test CIDEr from 88.7 to 97.1. With beam
search, the CE model gains much more from
the search, and the difference shrinks to 0.3 CIDEr, while BLEU-4 drops by 2.7
points. All SCST variants also narrow the output: the vocabulary falls from 419
words to between 212 and 243, and the share of unique captions from 50\% to
between 28\% and 33\%.

\textit{Pure versus mixed SCST.} Pure SCST reaches the highest CIDEr (98.6),
but 77.9\% of its test captions end in a function word. Most end in ``with a''
(64\% of all test captions), as in ``a plane is sitting on the runway with a''
or ``a kitchen with a sink and a'', where the mixed model writes ``a kitchen
with a sink and a window''. CIDEr does not penalize these unfinished endings.
With 20\% of the CE loss kept, no caption ends this way, at a
cost of 1.8 CIDEr. The mixed model also uses a larger vocabulary (241 versus
212 words) and scores slightly higher on BLEU-4
and ROUGE-L. The mixed objective does not stop the loss of diversity: unique
captions fall to 29\% with it and to 33\% with pure SCST.

\textit{Advantage normalization.} Without normalization, the policy-gradient
term is about ten times larger at the start of SCST, and validation CIDEr drops
from 88.1 to 81.2 in the first epoch, compared with 85.9 with normalization.
The model recovers, but after 10 epochs it remains 1.3 CIDEr below the
normalized run on the test split.

\textit{Optimizer reset.} Keeping the Adam state from the CE stage gives 97.6
CIDEr, 0.8 points above the run with a reset. With one run per setting, we do
not find a benefit from the reset.

\begin{table*}[t]
\caption{Training-objective ablation (Karpathy test, scores $\times100$). All SCST
variants start from the same CE checkpoint. Right block: caption statistics used
to measure reward hacking (Vocab: distinct words used; Uniq.: \% unique captions;
D-2: distinct-2; Rep.: \% captions with a repeated content word; Dangl.: \%
captions ending in a function word; Len.: mean length).}
\label{tab:loss}
\centering
\setlength{\tabcolsep}{4.5pt}
\begin{tabular}{@{}lccccc|cccccc@{}}
\toprule
Objective & B-1 & B-4 & M & R & C & Vocab & Uniq. & D-2 & Rep. & Dangl. & Len. \\
\midrule
CE only & 72.0 & 31.3 & 25.0 & 53.3 & 96.5 & 419 & 50.4 & 3.5 & 13.4 & 0.0 & 9.4 \\
Pure SCST ($\lambda=1$) & 73.7 & 28.2 & 23.9 & 52.1 & 98.6 & 212 & 32.7 & 1.7 & 1.3 & 77.9 & 9.7 \\
Mixed SCST ($\lambda=0.8$) & 72.7 & 28.6 & 24.0 & 52.8 & 96.8 & 241 & 29.3 & 1.6 & 2.5 & 0.0 & 9.1 \\
\quad w/o advantage normalization & 72.1 & 27.9 & 23.8 & 52.5 & 95.5 & 236 & 28.5 & 1.5 & 4.2 & 0.0 & 9.1 \\
\quad w/o optimizer reset & 72.8 & 28.8 & 24.2 & 53.0 & 97.6 & 243 & 31.3 & 1.7 & 5.2 & 0.0 & 9.1 \\
\bottomrule
\end{tabular}
\end{table*}

\subsection{Training Dynamics}
Fig.~\ref{fig:dynamics} shows validation CIDEr and the dangling-ending rate
over SCST training. Every variant loses CIDEr in the first epoch; the runs
with advantage normalization recover above the CE level by epoch 3, the run
without it only in epoch 4. Pure SCST starts producing dangling endings in
epoch 3, and they reach 76.4\% of validation captions by epoch 10, while the
mixed run stays at or below 0.04\%. In all variants, unique captions fall from 64\% to
between 32\% and 38\% within the first epoch.

\begin{figure}[t]
\centering
\figslot{figs/training_dynamics.pdf}{\columnwidth}{training dynamics}
\caption{Validation CIDEr (greedy decoding, top) and share of captions ending
in a function word (bottom) per SCST epoch, starting from the same CE
checkpoint (epoch 0).}
\label{fig:dynamics}
\end{figure}

\subsection{Effect of Decoding Heuristics}
\label{sec:decoding}
On the test split, the CE model scores 88.7 CIDEr with greedy decoding, 96.0
with plain beam search (beam 3), and 96.5 with our protocol (beam 7 with length
penalty, minimum length, and $n$-gram blocking). The mixed SCST model scores
97.1, 96.8, and 96.8. The heuristics therefore change CIDEr by at most 0.5
points, so the results do not depend on them. Beam search itself adds 7.3 to
7.8 CIDEr for the CE model, while for the mixed SCST model it lowers CIDEr
slightly (97.1 to 96.8).

\subsection{From-Scratch Run}
\label{sec:original}
Our first experiment trained an SE-ResNet-50 with a deep $3\times3$ stem from
scratch on MS-COCO, without ImageNet pretraining. Encoder and decoder were
trained together in three stages: CE with a frozen encoder, CE fine-tuning of
the deeper encoder stages, and mixed SCST. We kept only the epoch-87
checkpoint. In it, the stem and the first residual stage are still at their
random initialization, since they stayed frozen in every stage; stages 2 to 4
were learned. The SCST stage used a learning rate of $10^{-6}$, which changes
the model only slightly (Section~\ref{sec:training}). With our evaluation
protocol, this model reaches 22.1 BLEU-4 and 61.0 CIDEr (Table~\ref{tab:main}).
This result, obtained without pretraining, motivated the controlled comparison
of pretrained encoders in Table~\ref{tab:encoder}.

\subsection{Qualitative Analysis}
Fig.~\ref{fig:qualitative} shows captions for three randomly drawn test
images. On the snowboarding image, the CE and mixed models reproduce the
reference exactly, while pure SCST stops at ``down a snow''; its other two
captions end in ``with a''. All models also make content errors: every model
calls the Christmas figurines teddy bears or vases, and the CE model repeats
itself (``a teddy bear sitting next to a teddy bear'').

\begin{figure}[t]
\centering
\figslot{figs/qualitative.pdf}{\columnwidth}{example captions}
\caption{Captions for three Karpathy test images drawn at random, from the CE,
mixed SCST, and pure SCST models (SE-ResNet-50 encoder, beam search), with one
of the reference captions (GT).}
\label{fig:qualitative}
\end{figure}

\subsection{Limitations}
Our encoders are frozen; fine-tuning them could improve all models.
We evaluate SE and ResNet-D separately because no SE-ResNet-50-D pretrained with
the same recipe is available. Each configuration was trained once, so we did
not measure run-to-run variation; we treat differences of about one CIDEr point
or less, such as SE versus the baseline and the optimizer-reset ablation, as
inconclusive. Finally, the diversity and repetition statistics measure specific failure modes and are not
a complete measure of fluency.

% ===========================================================================
\section{Conclusion}
% ===========================================================================
Among frozen ImageNet-pretrained encoders, the ResNet-D modifications improved
CIDEr by 1.5 points over ResNet-50, while SE blocks gave no clear gain. Pure
SCST learned to end most captions with ``with a''; keeping 20\% of the
cross-entropy loss removed these endings for a cost of 1.8 CIDEr, and
normalizing the advantage softened the drop at the switch to reinforcement
learning. No SCST variant kept the caption diversity of the CE model, and with
beam search their CIDEr gain over CE was small. Future work includes
fine-tuning the encoder, combining SE blocks with the ResNet-D modifications
under a common pretraining recipe, multi-sample SCST
baselines~\cite{luo2020better}, and semantic rewards such as SPICE.

% ===========================================================================
% References
% ===========================================================================
\begin{thebibliography}{00}
\bibitem{vinyals2015show} O. Vinyals, A. Toshev, S. Bengio, and D. Erhan, ``Show and tell: A neural image caption generator,'' in \textit{Proc. IEEE Conf. Comput. Vis. Pattern Recognit. (CVPR)}, 2015, pp. 3156--3164.
\bibitem{xu2015show} K. Xu \textit{et al.}, ``Show, attend and tell: Neural image caption generation with visual attention,'' in \textit{Proc. Int. Conf. Mach. Learn. (ICML)}, 2015, pp. 2048--2057.
\bibitem{karpathy2015deep} A. Karpathy and L. Fei-Fei, ``Deep visual-semantic alignments for generating image descriptions,'' in \textit{Proc. IEEE Conf. Comput. Vis. Pattern Recognit. (CVPR)}, 2015, pp. 3128--3137.
\bibitem{mao2015deep} J. Mao, W. Xu, Y. Yang, J. Wang, Z. Huang, and A. Yuille, ``Deep captioning with multimodal recurrent neural networks (m-RNN),'' in \textit{Proc. Int. Conf. Learn. Represent. (ICLR)}, 2015.
\bibitem{ranzato2016sequence} M. Ranzato, S. Chopra, M. Auli, and W. Zaremba, ``Sequence level training with recurrent neural networks,'' in \textit{Proc. Int. Conf. Learn. Represent. (ICLR)}, 2016.
\bibitem{bengio2015scheduled} S. Bengio, O. Vinyals, N. Jaitly, and N. Shazeer, ``Scheduled sampling for sequence prediction with recurrent neural networks,'' in \textit{Adv. Neural Inf. Process. Syst. (NeurIPS)}, 2015, pp. 1171--1179.
\bibitem{rennie2017self} S. J. Rennie, E. Marcheret, Y. Mroueh, J. Ross, and V. Goel, ``Self-critical sequence training for image captioning,'' in \textit{Proc. IEEE Conf. Comput. Vis. Pattern Recognit. (CVPR)}, 2017, pp. 7008--7024.
\bibitem{papineni2002bleu} K. Papineni, S. Roukos, T. Ward, and W.-J. Zhu, ``BLEU: A method for automatic evaluation of machine translation,'' in \textit{Proc. Annu. Meeting Assoc. Comput. Linguistics (ACL)}, 2002, pp. 311--318.
\bibitem{vedantam2015cider} R. Vedantam, C. L. Zitnick, and D. Parikh, ``CIDEr: Consensus-based image description evaluation,'' in \textit{Proc. IEEE Conf. Comput. Vis. Pattern Recognit. (CVPR)}, 2015, pp. 4566--4575.
\bibitem{banerjee2005meteor} S. Banerjee and A. Lavie, ``METEOR: An automatic metric for MT evaluation with improved correlation with human judgments,'' in \textit{Proc. ACL Workshop Intrinsic Extrinsic Eval. Measures Mach. Transl. Summarization}, 2005, pp. 65--72.
\bibitem{lin2004rouge} C.-Y. Lin, ``ROUGE: A package for automatic evaluation of summaries,'' in \textit{Text Summarization Branches Out}, 2004, pp. 74--81.
\bibitem{lin2014microsoft} T.-Y. Lin \textit{et al.}, ``Microsoft COCO: Common objects in context,'' in \textit{Proc. Eur. Conf. Comput. Vis. (ECCV)}, 2014, pp. 740--755.
\bibitem{chen2015microsoft} X. Chen \textit{et al.}, ``Microsoft COCO captions: Data collection and evaluation server,'' \textit{arXiv preprint arXiv:1504.00325}, 2015.
\bibitem{he2016deep} K. He, X. Zhang, S. Ren, and J. Sun, ``Deep residual learning for image recognition,'' in \textit{Proc. IEEE Conf. Comput. Vis. Pattern Recognit. (CVPR)}, 2016, pp. 770--778.
\bibitem{he2019bag} T. He, Z. Zhang, H. Zhang, Z. Zhang, J. Xie, and M. Li, ``Bag of tricks for image classification with convolutional neural networks,'' in \textit{Proc. IEEE Conf. Comput. Vis. Pattern Recognit. (CVPR)}, 2019, pp. 558--567.
\bibitem{hu2018squeeze} J. Hu, L. Shen, and G. Sun, ``Squeeze-and-excitation networks,'' in \textit{Proc. IEEE Conf. Comput. Vis. Pattern Recognit. (CVPR)}, 2018, pp. 7132--7141.
\bibitem{wightman2021resnet} R. Wightman, H. Touvron, and H. J\'egou, ``ResNet strikes back: An improved training procedure in timm,'' \textit{arXiv preprint arXiv:2110.00476}, 2021.
\bibitem{anderson2018bottom} P. Anderson \textit{et al.}, ``Bottom-up and top-down attention for image captioning and visual question answering,'' in \textit{Proc. IEEE Conf. Comput. Vis. Pattern Recognit. (CVPR)}, 2018, pp. 6077--6086.
\bibitem{huang2019attention} L. Huang, W. Wang, J. Chen, and X.-Y. Wei, ``Attention on attention for image captioning,'' in \textit{Proc. IEEE/CVF Int. Conf. Comput. Vis. (ICCV)}, 2019, pp. 4634--4643.
\bibitem{cornia2020meshed} M. Cornia, M. Stefanini, L. Baraldi, and R. Cucchiara, ``Meshed-memory transformer for image captioning,'' in \textit{Proc. IEEE/CVF Conf. Comput. Vis. Pattern Recognit. (CVPR)}, 2020, pp. 10578--10587.
\bibitem{paulus2018deep} R. Paulus, C. Xiong, and R. Socher, ``A deep reinforced model for abstractive summarization,'' in \textit{Proc. Int. Conf. Learn. Represent. (ICLR)}, 2018.
\bibitem{luo2020better} R. Luo, ``A better variant of self-critical sequence training,'' \textit{arXiv preprint arXiv:2003.09971}, 2020.
\end{thebibliography}

\end{document}
